\section{Height moments and planar nesting}\label{sec:projection}

The cap coefficient has growth $N^{1/6+o(1)}$ at winding fugacity $2$.
To extract a planar nesting exponent, we must control the height of
polygons before projecting the cylinder. Arbitrary fixed twist derivatives
provide the required height moments.

\subsection{Uniform bounds for every fixed twist derivative}
\label{subsec:v-higher-twist-jets}

The first winding derivative controls the total mass of one winding
polygon.  We now bound every fixed derivative.  These bounds will control
the heights of the polygons before passage from the cylinder to the plane.
The analyticity radius in the winding fugacity may depend on the
circumference; no uniform radius will be needed.

We use the following notation from Section~\ref{sec:cylinder}:
\[
  l=\frac{\pi}{8},\qquad
 \alpha_{\mathrm s}=\frac{\pi}{6 l}=\frac43,\qquad
 \beta_{\mathrm s}=\alpha_{\mathrm s}^{-1}=\frac34,
 \qquad S_a=\{s\in\mathbb C:|\Re s|<a\},
\]
\[
 h(s)=\frac{\cos s-\cos3 l}{\cos s+\cos l},\qquad
 B(s)=\frac{\cos(\alpha_{\mathrm s}s)-d_0}
             {\cos(\alpha_{\mathrm s}s)+d_0},\qquad
 d_0=\cos(\alpha_{\mathrm s} l)=\frac{\sqrt3}{2}.
\]
An even holomorphic function on $S_{3 l}$ is a holomorphic function
of $z=B(s)$ on the disk.  The point $s= l$ maps to $z=0$, and
$B'( l)\ne0$.  Put
\[
 (\mathcal C_g   U)(s)=\int_{\mathbb R}g(t)U(s+it)\,dt,
 \qquad
 g(t)=\frac1{2\pi}
 \left(\frac{2^{-1/2}}{\cosh t-2^{-1/2}}-\frac1{\cosh t}\right).
\]
Thus $p:=\|g\|_1<1/2$ and $\int g=1/4$.

For precision, here are all the operator facts from the strip calculation in Section~\ref{sec:cylinder} that enter the argument.  If $j=f\circ h$ with $\deg f\le N$,
then, on $S_{3 l}$,
\begin{equation}\label{eq:v-jet-strip-representation}
 j(s-2 l)+j(s+2 l)-j(s)
       =(I-\mathcal C_g  )q(B(s))+\lambda,
 \qquad q(z)=\sum_{m=1}^N q_mz^m.
\end{equation}
Here $q$ has no constant term.  With the inverse branch $s(0)= l$, define
\[
 M_{nm}=[z^n](\mathcal C_g   B^m)(s(z))\quad(n,m\ge1),
 \qquad M_{0m}=(\mathcal C_g  B^m)( l).
\]
The matrix $M$ is real and self-adjoint on the weighted space
$\ell^2(n)$ with squared norm $\sum_{n\ge1}n|v_n|^2$, and its operator
norm is at most $p$.  For one constant $C$ independent of the indices,
\begin{equation}\label{eq:v-jet-matrix-inputs}
 |M_{nm}|\le\frac Cn
 \min\left\{\left(\frac nm\right)^{\beta_{\mathrm s}},
             \left(\frac mn\right)^{\beta_{\mathrm s}}\right\},
 \qquad
 \sum_{n\ge1}|M_{nm}|\longrightarrow p,
 \qquad M_{0m}\longrightarrow0.
\end{equation}
These are operator inputs, not assumptions on the higher derivatives.
In particular, we will not assume that $M$ is a contraction on $\ell^1$.

\begin{lemma}[Higher twist derivatives]\label{lem:v-higher-twist-jets}
For each $N\ge1$, suppose the finite-row eigenvalue has an expansion
\[
 F_N(s,\nu)=1+\sum_{k\ge1}\nu^k j_{N,k}(s),
 \qquad j_{N,k}(s)=f_{N,k}(h(s)),\qquad \deg f_{N,k}\le N,
\]
analytic near $\nu=0$, and satisfies $f_{N,k}(1)=\mathbf1_{k=1}$.
Suppose also that every coefficient of
\[
 F_N(s-2 l,\nu)F_N(s+2 l,\nu)-F_N(s,\nu)
\]
vanishes to order at least $N$ at $s= l$.
Assume the strip-operator facts
\eqref{eq:v-jet-strip-representation}--\eqref{eq:v-jet-matrix-inputs}.
Set, with the index $N$ suppressed,
\[
 \begin{split}
 H_k(s)&=j_k(s-2 l)+j_k(s+2 l)-j_k(s),\\
 P_k(s)&=\sum_{i=1}^{k-1}j_i(s-2 l)j_{k-i}(s+2 l),
 \qquad G_k=H_k+P_k.
 \end{split}
\]
For each fixed $k\ge1$ there is a constant $C_k$, independent of $N$,
such that
\begin{equation}\label{eq:v-jet-full-strip-conclusion}
 \sup_{s\in S_{5 l}}|j_k(s)|\le C_k,
 \qquad |G_k(s)|\le C_k|B(s)|^N\quad(s\in S_{3 l}).
\end{equation}
For every $0<\delta<5 l$ there is also a constant $C_{k,\delta}$,
independent of $N$, such that
\begin{equation}\label{eq:v-jet-sharp-conclusion}
 |j_k(s)|\le C_{k,\delta}
       \min\left\{1,\frac{e^{\alpha_{\mathrm s}|\Im s|}}N\right\}
 \qquad (|\Re s|\le5 l-\delta).
\end{equation}
Finally,
\begin{equation}\label{eq:v-jet-endpoint-conclusion}
 |f_{N,k}'(1)|\le C_k N^{\beta_{\mathrm s}},\qquad
 |[\nu^k]\log F_N(- l,\nu)|\le\frac{C_k}{N}.
\end{equation}
\end{lemma}

\begin{proof}
All divisions in the variable $\nu$ below are divisions of formal power
series with constant coefficient one.  At a fixed order they use only
finitely many coefficients.  Thus they require neither a uniform
analyticity radius in $\nu$ nor a lower bound on $|F_N|$ away from zero.

\paragraph{A uniform inverse for the coefficient equations.}
Let $M^{(N)}=(M_{nm})_{1\le n,m<N}$.  The Hilbert-space contraction
implies that $I-M^{(N)}$ is invertible.  In fact
\begin{equation}\label{eq:v-jet-uniform-lone-inverse}
 \sup_{N\ge1}\|(I-M^{(N)})^{-1}\|_{\ell^1\to\ell^1}<\infty.
\end{equation}
To prove this, note from \eqref{eq:v-jet-matrix-inputs} that the column
sums $d_m:=\sum_n|M_{nm}|$ are uniformly bounded, that $d_m\to p<1$,
and that every fixed row tends to zero as $m\to\infty$.  If
\eqref{eq:v-jet-uniform-lone-inverse} failed, there would be
$N_r\to\infty$ and vectors $v^{(r)}$, extended by zero outside
$1,\ldots,N_r-1$, with
\[
 \|v^{(r)}\|_1=1,\qquad
 \|v^{(r)}-\mathbf1_{\{n<N_r\}}Mv^{(r)}\|_1\longrightarrow0.
\]
Pass to a coordinatewise limit $v\in\ell^1$.  The vanishing tails of
each row imply $v=Mv$.  Moreover, the entry bound gives
$|v_n|\le C/n$, and hence
\[
 \sum_n n|v_n|^2\le C\sum_n|v_n|<\infty.
\]
The strict contraction on $\ell^2(n)$ therefore forces $v=0$.
Choose $J$ so large that $d_m\le p+\varepsilon<1$ for $m>J$.
The approximate equation then gives the contradiction
\[
 1\le\sum_{m\le J}d_m|v_m^{(r)}|
       +(p+\varepsilon)\|v^{(r)}\|_1+o(1)
   \longrightarrow p+\varepsilon<1.
\]

We also need the normalization of the one-dimensional homogeneous
solution.  Let
\[
 Q_N(z)=\sum_{m=1}^N c_mz^m,\qquad c_N=1,\qquad
 c_n=\sum_{m=1}^N M_{nm}c_m\quad(1\le n<N).
\]
Equation \eqref{eq:v-jet-uniform-lone-inverse} and the bounded column
sums give $\sup_N\sum_m|c_m|<\infty$.  The same coordinatewise-limit
argument shows $c_m\to0$ for every fixed $m$ as $N\to\infty$.  Using
the high-column bound once more yields
\begin{equation}\label{eq:v-jet-homogeneous-lower-coefficients}
 \limsup_{N\to\infty}\sum_{m<N}|c_m|
       \le\frac p{1-p}<1.
\end{equation}
For a polynomial $q$ with no constant term, put
\[
 (\mathcal Tq)(s)=(I-\mathcal C_g  )q(B(s))
                         +(\mathcal C_g  q(B))( l),
\]
\begin{equation}\label{eq:v-jet-normalization-functional}
 \mathcal E(q):=(\mathcal Tq)(i\infty)
       =\frac34q(1)+\sum_mM_{0m}q_m.
\end{equation}
Since $M_{0m}\to0$, the last sum tends to zero for $q=Q_N$.
The coefficients $c_m$ are real.  Consequently
\begin{equation}\label{eq:v-jet-positive-endpoint}
 \liminf_{N\to\infty}\mathcal E(Q_N)
       \ge\frac34\frac{1-2p}{1-p}>0.
\end{equation}
For completeness, $\mathcal E(Q_N)$ is nonzero at every remaining
finite $N$ as well.  Indeed the existing first jet has $P_1=0$;
its order-$N$ zero makes the polynomial in
\eqref{eq:v-jet-strip-representation} a multiple of $Q_N$, while its
endpoint is $H_1(i\infty)=1$.  Thus $\mathcal E(Q_N)=0$ is impossible.
Together with \eqref{eq:v-jet-positive-endpoint}, this gives a uniform
positive lower bound on $|\mathcal E(Q_N)|$.

\paragraph{The forcing coefficients.}
We proceed by induction on $k$.  Assume the two strip bounds in the
lemma for every lower order.  When $k=1$ there are no lower-order
terms and the forcing below is zero.  For $k\ge2$, let
\[
 b_m=[z^m]P_k(s(z)).
\]
The function $P_k$ is even: replacing $s$ by $-s$ and $i$ by $k-i$
leaves its sum unchanged.  Hence it descends holomorphically to the
disk.  We claim
\begin{equation}\label{eq:v-jet-forcing-lone}
 \sum_{m=0}^{N-1}|b_m|\le C_k.
\end{equation}

Here is a contour proof, including the cancellation across its slit.
Temporarily write
\[
 A=F_N(s-2 l,\nu),\quad C=F_N(s,\nu),\quad
 D=F_N(s+2 l,\nu),\quad E=AD-C,
 \qquad a_i=[\nu^i]A.
\]
Define the replacement
\[
 R_k(s)=\sum_{i=1}^{k-1}a_i[\nu^{k-i}](C/A).
\]
The exact identity $D-C/A=E/A$ gives
\begin{equation}\label{eq:v-jet-replacement-identity}
 R_k-P_k=-\sum_{i=1}^{k-1}a_i[\nu^{k-i}](E/A).
\end{equation}
Only defect coefficients $G_r$ with $r<k$ occur on the right.  Each is
$O(B^N)$ near $s= l$.  Therefore
\begin{equation}\label{eq:v-jet-replacement-coefficients}
 [z^m]R_k(s(z))=b_m\qquad(0\le m<N).
\end{equation}

Put $b_*:=B(0)=(1-d_0)/(1+d_0)$.  The inverse branch continued from
$s(0)= l$ is
\[
 s(z)=\alpha_{\mathrm s}^{-1}
       \arccos\left(d_0\frac{1+z}{1-z}\right).
\]
Fix $0<\varepsilon_s< l/4$.  There are fixed $\eta>0$ and
$\theta\in(\pi/2,\pi/2+\alpha_{\mathrm s}\varepsilon_s)$ such that
this branch is holomorphic on
\[
 \mathcal D=\{z:|z|<1+\eta,\quad
                    |\arg(1-z)|<\theta\}\setminus[b_*,1),
\]
and satisfies
\begin{equation}\label{eq:v-jet-inverse-domain}
 0<\Re s(z)<3 l+\varepsilon_s,
 \qquad e^{\alpha_{\mathrm s}|\Im s(z)|}\le\frac C{|1-z|}.
\end{equation}
These facts follow directly from the inverse formula: away from $1$
one uses compactness and the disk mapping, while near $1$ the logarithmic
inverse has real part at most
$\alpha_{\mathrm s}^{-1}|\arg(1-z)|+o(1)$.
The other branch point, $(1+d_0)/(1-d_0)$, is excluded by taking $\eta$
small.  On the two banks of $b_*<z<1$, the boundary values of the
inverse are $\pm iY(z)$, where
\[
 Y(z)=\alpha_{\mathrm s}^{-1}
       \operatorname{arcosh}\left(d_0\frac{1+z}{1-z}\right),
 \qquad e^{\alpha_{\mathrm s}Y(z)}\asymp(1-z)^{-1}.
\]

Throughout \eqref{eq:v-jet-inverse-domain}, the arguments $s$ and
$s-2 l$ lie in a fixed closed inner substrip of $S_{5 l}$.
Every lower-order jet at either argument is bounded by
$C_k t_N(s)$, where
\[
 t_N(s)=\min\{1,e^{\alpha_{\mathrm s}|\Im s|}/N\}.
\]
The fixed-order coefficients of $A^{-1}$ are bounded, and those of
$C/A-1=(C-A)/A$ are $O_k(t_N)$.  Thus the replacement retains two small
jet factors:
\begin{equation}\label{eq:v-jet-replacement-size}
 |R_k(s(z))|\le C_k\min\{1,(N|1-z|)^{-2}\}
       \qquad(z\in\mathcal D).
\end{equation}
On either bank, \eqref{eq:v-jet-replacement-identity} retains one such
factor and one lower-order defect, giving
\[
 |R_k(\pm iY(z))-P_k(\pm iY(z))|
       \le C_k z^N\min\{1,(N(1-z))^{-1}\}.
\]
The two values of $P_k$ agree by evenness.  The jump of the replacement
therefore satisfies
\begin{equation}\label{eq:v-jet-bank-jump}
 |R_{k,+}(z)-R_{k,-}(z)|
       \le\frac{C_k z^N}{N(1-z)}\qquad(b_*<z<1).
\end{equation}

Choose fixed $1<R<1+\eta$ and $\pi/2<\theta'<\theta$.
For all sufficiently large $N$, deform the coefficient contour to the
appropriate outer arc $|z|=R$, the two segments
$z=1-r e^{\pm i\theta'}$ down to $r=1/N$, the two joining arcs of
$|z-1|=1/N$, and the oppositely oriented slit banks from $1-1/N$ to
$b_*$.  A small circle about $b_*$ has vanishing contribution: the
replacement is bounded there.  The joining arcs are interrupted at
$1-1/N$ by the two banks.  All contours may first be taken strictly
inside $\mathcal D$ and then passed to their boundary values.

For $m<N$, the non-bank parts of the Cauchy integral for
$R_k(s(z))z^{-m-1}$ cost $O_k(N^{-1})$.  On the outer arc the stronger
bound $O_k(N^{-2})$ holds.  On the two straight segments, $|z|\ge1$
because $\cos\theta'<0$, and
\[
 N^{-2}\int_{1/N}^{O(1)}r^{-2}\,dr=O(N^{-1}).
\]
On the joining arcs, $|z|\ge1-1/N$, so
$|z|^{-m-1}\le(1-1/N)^{-N}\le4$ for $N\ge2$; their length is
$O(N^{-1})$ and the replacement is bounded.  Summing these estimates
over $0\le m<N$ costs $O_k(1)$.
The banks must be combined before taking absolute values.  By
\eqref{eq:v-jet-bank-jump}, their total contribution to the sum of
coefficient bounds is at most
\[
 \frac{C_k}{N}\int_{b_*}^{1-1/N}
       \frac{\sum_{m=0}^{N-1}z^{N-m-1}}{1-z}\,dz
 \le\frac{C_k}{N}\int_{b_*}^{1-1/N}\frac{dz}{(1-z)^2}
 \le C_k.
\]
Together with \eqref{eq:v-jet-replacement-coefficients}, this proves
\eqref{eq:v-jet-forcing-lone} for all sufficiently large $N$.
The finitely many smaller $N$ use an ordinary fixed circle
$|z|=r_*<b_*$ and the preceding-jet bounds there.  This also covers
$N=1$, for which the indentation $|z-1|=1/N$ must not be used.

\paragraph{Solving and normalizing the current jet.}
Apply \eqref{eq:v-jet-strip-representation} to $j_k$.  Since $G_k$
vanishes to order $N$ at $z=0$, its lower coefficients give
\begin{equation}\label{eq:v-jet-forced-equations}
 q_n-(Mq)_n=-b_n\quad(1\le n<N),\qquad
 \lambda=(\mathcal C_g  q(B))( l)-b_0.
\end{equation}
Let $u_N=0$ and
$u_{<N}=-(I-M^{(N)})^{-1}b_{<N}$.  By
\eqref{eq:v-jet-uniform-lone-inverse} and
\eqref{eq:v-jet-forcing-lone}, $\|u\|_1\le C_k$.
Every solution of the first equations in
\eqref{eq:v-jet-forced-equations} has the form $q=u+tQ_N$.
The endpoint condition is on $H_k$, not on $G_k$:
\[
 H_k(i\infty)=\mathbf1_{k=1},\qquad
 H_k=\mathcal Tq-b_0.
\]
Consequently
\[
 t=\frac{\mathbf1_{k=1}+b_0-\mathcal E(u)}{\mathcal E(Q_N)}.
\]
The numerator is bounded, since
$|\mathcal E(u)|\le(3/4+p)\|u\|_1$, and the denominator is bounded
away from zero.  We have proved
\begin{equation}\label{eq:v-jet-q-bound}
 \|q\|_1+|\lambda|\le C_k,
 \qquad \sup_{S_{3 l}}|H_k|\le C_k.
\end{equation}
Notice that $P_k(i\infty)=\mathbf1_{k=2}$; replacing the stated
endpoint of $H_k$ by an endpoint of $G_k$ would give the wrong
normalization.

\paragraph{Inversion and bounds up to the strip edges.}
For each fixed $N,k$, the polynomial form gives
$j_k(x+iy)-\mathbf1_{k=1}=O_{N,k,\eta}(e^{-|y|})$ on
$|x|\le7 l-\eta$.  This endpoint behavior justifies Fourier
transformation along vertical lines and horizontal contour displacement.
The linear shift multiplier is $2\cosh(2 l\xi)-1$, whose inverse is
$\cosh( l\xi)/\cosh(3 l\xi)$.  Its Fourier kernel is
\begin{equation}\label{eq:v-jet-inverse-kernel}
 \mathcal K(t)=
 \frac{\operatorname{sech}(\alpha_{\mathrm s}t+i\pi/6)
       +\operatorname{sech}(\alpha_{\mathrm s}t-i\pi/6)}{12 l}
 =\frac{\sqrt3\cosh(\alpha_{\mathrm s}t)}
        {12 l(\sinh^2(\alpha_{\mathrm s}t)+3/4)}.
\end{equation}
On the real axis this kernel is positive, has integral one, and obeys
\[
 \mathcal K(t)\le\frac{\sqrt3}{6 l}
                           e^{-\alpha_{\mathrm s}|t|}.
\]
Thus
\begin{equation}\label{eq:v-jet-inverse-formula}
 j_k(x+iy)=\int_{\mathbb R}\mathcal K(t)
                       H_k(x+i(y-t))\,dt\qquad(|x|<3 l).
\end{equation}
The endpoint condition is essential here: analyticity in a strip alone
would not exclude exponentially growing solutions of the homogeneous
shift equation.

Equation \eqref{eq:v-jet-inverse-formula} bounds $j_k$ on the whole
open strip $S_{3 l}$.  If $3 l\le\Re s<5 l$, use the exact
rearrangement
\begin{equation}\label{eq:v-jet-edge-rearrangement}
 j_k(s)=H_k(s-2 l)+j_k(s-2 l)-j_k(s-4 l).
\end{equation}
All three arguments on the right belong to $S_{3 l}$.  The opposite
edge follows by evenness.  In particular
\[
 \sup_{S_{5 l}}|j_k|\le3\sup_{S_{3 l}}|H_k|\le C_k.
\]
This is a bound on the entire open strip, uniform up to its edges.
It is obtained before any new bound on $G_k$ is used.

The lower-order full-strip bounds now make $P_k$ bounded on
$S_{3 l}$.  Hence $G_k$ is bounded there as well.  It is even and
has an order-$N$ zero at $B=0$.  The disk maximum principle gives
\[
 |G_k(s)|\le C_k|B(s)|^N\qquad(s\in S_{3 l}).
\]
For every $\eta>0$, the explicit formula for $B$ gives
\[
 |B(s)|\le\exp(-c_\eta e^{-\alpha_{\mathrm s}|\Im s|})
       \quad(|\Re s|\le3 l-\eta).
\]
Since $e^{-cu}\le C_c\min\{1,u^{-2}\}$, the last bound on $G_k$
and the inductive sharp bounds for the two factors of $P_k$ imply
\begin{equation}\label{eq:v-jet-sharp-forcing}
 |H_k(s)|\le C_{k,\eta}
       \min\left\{1,\frac{e^{2\alpha_{\mathrm s}|\Im s|}}{N^2}\right\}
       \quad(|\Re s|\le3 l-\eta).
\end{equation}
Convolution improves the squared scale to the desired scale.  Indeed,
with $a=e^{\alpha_{\mathrm s}|y|}/N$, the case $a\ge1$ uses
$\int\mathcal K=1$.  When $a<1$, split at
$T=\alpha_{\mathrm s}^{-1}\log(1/a)$ and use
\[
 \int_{\mathbb R}e^{-\alpha_{\mathrm s}|t|}
       \min\{1,a^2e^{2\alpha_{\mathrm s}|t|}\}\,dt
       =\frac{4a-2a^2}{\alpha_{\mathrm s}}
       \le\frac{4a}{\alpha_{\mathrm s}}.
\]
Equations \eqref{eq:v-jet-inverse-formula} and
\eqref{eq:v-jet-sharp-forcing} therefore prove
\eqref{eq:v-jet-sharp-conclusion} on each closed inner substrip of
$S_{3 l}$.  To reach the asserted width, take
$\eta=\min\{\delta, l\}$.  For
$3 l-\eta\le\Re s\le5 l-\eta$, all arguments on the right
of \eqref{eq:v-jet-edge-rearrangement} lie in
$|\Re s|\le3 l-\eta$.  Apply the sharp bounds there, and use
$\min\{1,a^2\}\le\min\{1,a\}$.  The negative edge is identical.
This proves \eqref{eq:v-jet-sharp-conclusion} and closes the induction.
For $k=1$, the argument starts with $P_1=0$ and no contour forcing.

\paragraph{The derivative at $h=1$.}
No sign of a higher jet is needed for this step.  Define
\[
 \kappa=2(\cos l+\cos3 l)(2\cos2 l-1)>0,
 \qquad \gamma_g=\frac{1-2^{-1/2}}\pi>0.
\]
For fixed $N,k$, expansion of the rational function $h$ at imaginary
infinity gives
\begin{equation}\label{eq:v-jet-endpoint-first-limit}
 \lim_{y\to\infty}e^y
       (\mathbf1_{k=1}-H_k(iy))=\kappa f_{N,k}'(1).
\end{equation}
For the polynomial $q$ found above, put
\[
 D_q(u)=q(1)-q(B(iu)),\qquad
 I_m=\int_{\mathbb R}e^u(1-B(iu)^m)\,du.
\]
The constant $\lambda$ cancels in the endpoint difference, so
\[
 \mathbf1_{k=1}-H_k(iy)
       =D_q(y)-\int_{\mathbb R}g(t)D_q(y+t)\,dt.
\]
For each fixed polynomial, $D_q(u)=O(e^{-\alpha_{\mathrm s}|u|})$.
Also $g(v)\sim-\gamma_g e^{-|v|}$ and
$e^y|g(u-y)|\le Ce^u$.  Because $\alpha_{\mathrm s}>1$, dominated
convergence gives
\[
 \lim_{y\to\infty}e^y
       (\mathbf1_{k=1}-H_k(iy))
       =\gamma_g\sum_{m=1}^N q_m I_m.
\]
Here $0<B(iu)<1$ and
\[
 0\le1-B(iu)^m\le\min\{1,Cm e^{-\alpha_{\mathrm s}|u|}\}.
\]
Splitting the positive half-line at
$\alpha_{\mathrm s}^{-1}\log(Cm)$ and bounding the negative half-line
directly gives $0\le I_m\le C m^{\beta_{\mathrm s}}$.  Therefore
\[
 f_{N,k}'(1)=\frac{\gamma_g}{\kappa}\sum_{m=1}^Nq_mI_m,
 \qquad |f_{N,k}'(1)|\le C_kN^{\beta_{\mathrm s}},
\]
by \eqref{eq:v-jet-q-bound}.  The endpoint limit was taken for fixed
$N,k$; no interchange with an $N\to\infty$ limit has been used.
Finally, \eqref{eq:v-jet-sharp-conclusion} gives
$j_{N,r}(- l)=O_r(N^{-1})$ for every fixed $r$.  The coefficient of
$\nu^k$ in $\log F_N(- l,\nu)$ is a finite sum of products of these
jets, each with at least one factor.  It is consequently $O_k(N^{-1})$.
\end{proof}

\paragraph{Application to the annular transfer.}
Normalize the simple eigenvector near winding fugacity $\nu=0$ by its
empty coordinate, and write its row eigenvalue as $F_N(s,\nu)$.
The spectral gap at the physical row and commutation give an analytic
choice for each fixed $N$.  Empty-row evaluation gives
\[
 F_N(s,\nu)=1+\nu\left(h^N+
       \sum C_0^r(1-h)^r h^{\sum_a s_a}
                    \Psi_\nu(\text{arches over the gaps})\right),
 \qquad C_0=2\cos3 l,
\]
where the sum is over interval arrangements on the labelled ring.
Any closed loop meeting this row must be winding; the gap-pairing
argument therefore applies as for the first derivative.  This formula
has degree at most $N$ in $h$ and gives
$\lim_{y\to\infty}F_N(iy,\nu)=1+\nu$, or equivalently
$f_{N,k}(1)=\mathbf1_{k=1}$.
The auxiliary fusion relation, closed against these same caps before
differentiation, gives the order-$N$ zero of the product defect at
$s=- l$, and evenness gives it at $s= l$.
All hypotheses of Lemma~\ref{lem:v-higher-twist-jets} are thus the
already established finite-row and strip identities.  Its constants are
uniform in $N$ for each fixed derivative order, which is the uniformity
needed in the subsequent polygon-height estimates.


\subsection{From derivatives to polygon height moments}

We now use the derivative bounds in positive path sums at the physical row
$t=l$. Throughout this subsection a polygon $P$ has \emph{bare weight}
$w(P)=x^{|P|}$, with no loop fugacity included. An open path still has
$w(\gamma)=x^{m(\gamma)}$. Thus a winding polygon in the gas with fugacity
$\nu$ has activity $\nu w(P)$; in particular, $P_N$ is the bare mass from
Proposition~\ref{prop:cylinder}, not the mass at fugacity $2$.

For a path or polygon let $r_-$ and $r_+$ be its lowest and highest
occupied rhombus rows. Write
\[
 \operatorname{top}(P)=r_+(P),\qquad
 \operatorname{ht}(P)=r_+(P)-r_-(P)+1.
\]
For an arc below the cut at row $0$, put
$\operatorname{dep}(\gamma)=-r_-(\gamma)$. Changing between row cuts and
centre heights costs only an absolute number of rows. Constants denoted
by $C_{\rm lat}$ below allow for this fixed lattice convention, and
$u_+=\max\{u,0\}$.

\paragraph{The winding pair coefficient.}
Let $\mathscr W_{N,M}$ be the finite set of simple winding polygons whose
occupied rows are among $-M,\ldots,-1$. Two polygons are incompatible if
they share a dual vertex; two copies of the same polygon are also declared
incompatible. Define the finite gas partition
\[
 \Xi_{N,M}(\nu)=
 \sum_{\substack{\mathcal B\subset\mathscr W_{N,M}\\
                  \mathcal B\text{ pairwise compatible}}}
       \prod_{\eta\in\mathcal B}\nu w(\eta).
\]
The empty set contributes one. Expanding its logarithm at $\nu=0$ gives
the exact finite identity
\begin{equation}\label{eq:winding-pair-finite}
 -2[\nu^2]\log\Xi_{N,M}(\nu)
   =\sum_{\gamma,\eta\in\mathscr W_{N,M}}
       w(\gamma)w(\eta)\mathbf1_{\{\gamma\not\sim\eta\}}.
\end{equation}
The sum is ordered and includes repetitions. In transfer notation,
$\Xi_{N,M}(\nu)=\langle0|(\mathsf T_\nu(l))^M|0\rangle$. For each fixed $N$, the
analytic gap at $\nu=0$ gives a nonzero analytic leading coefficient
times $F_N(-l,\nu)^M$, with an analytic exponentially smaller remainder.
Consequently
\[
 \lim_{M\to\infty}\frac1M[\nu^2]\log\Xi_{N,M}(\nu)
       =[\nu^2]\log F_N(-l,\nu).
\]
This limit is taken for fixed $N$; it requires no uniform analyticity
radius in $\nu$.

Let $\mathscr W_N$ be the winding polygons at all row positions on the
infinite cylinder and let $\mathscr W_N^0$ be those with top row $0$.
Horizontal translates are distinct members when their vertex sets are
distinct, and $P_N=\sum_{\eta\in\mathscr W_N^0}w(\eta)$. For an anchored
$\gamma\in\mathscr W_N^0$ and a placed
$\eta\in\mathscr W_N$, let $r(\gamma,\eta)$ be the number of rows from the
lowest to the highest occupied row of their union. Translating the pair
vertically puts it in the $M$-row slab in exactly
$(M-r(\gamma,\eta)+1)_+$ ways. Hence the right side of
\eqref{eq:winding-pair-finite}, divided by $M$, is
\[
 \sum_{\gamma\in\mathscr W_N^0}\sum_{\eta\in\mathscr W_N}
 w(\gamma)w(\eta)\mathbf1_{\{\gamma\not\sim\eta\}}
       \frac{(M-r(\gamma,\eta)+1)_+}{M}.
\]
Each placement fraction increases to one. Monotone convergence and
Lemma~\ref{lem:v-higher-twist-jets} therefore give the finite mass
\begin{equation}\label{eq:winding-pair-mass}
 \begin{split}
 Q_N&:=\sum_{\gamma\in\mathscr W_N^0}w(\gamma)
       \sum_{\eta\in\mathscr W_N}
          w(\eta)\mathbf1_{\{\gamma\not\sim\eta\}}\\
   &=-2[\nu^2]\log F_N(-l,\nu)\le \frac{C}{N}.
 \end{split}
\end{equation}

We relate this incompatibility mass to height. A simple winding curve
separates the cylinder into an upper and a lower component. Write
$u_\eta(z)$ for the indicator that a lattice centre $z$ lies strictly in
the upper component of $\eta$. Let $\tau_{a,b}$ translate by
$a+b e^{i\pi/3}$, where $a\in\mathbb Z/N\mathbb Z$ and $b\in\mathbb Z$.
On a fixed $\gamma\in\mathscr W_N^0$, choose centres $z_+,z_-$ of the same
triangle type within $C_{\rm lat}$ rows of its two extreme rows. Their
row-height difference $\delta_\gamma$ satisfies
$\delta_\gamma\ge\operatorname{ht}(\gamma)-C_{\rm lat}$. Since the two
centres differ by one of the translations $\tau_{a,b}$, reindexing $a$ and
telescoping in $b$ gives, for every winding shape $\eta$,
\begin{equation}\label{eq:winding-indicator-telescope}
 \begin{split}
 N\delta_\gamma
 &=\sum_{a=0}^{N-1}\sum_{b\in\mathbb Z}
       \bigl(u_{\tau_{a,b}\eta}(z_+)
             -u_{\tau_{a,b}\eta}(z_-)\bigr)\\
 &\le\sum_{a=0}^{N-1}\sum_{b\in\mathbb Z}
       \mathbf1_{\{\gamma\not\sim\tau_{a,b}\eta\}}.
 \end{split}
\end{equation}
The differences vanish for all sufficiently high or low translates, so
the first sum is finite. For the inequality, if $\eta$ avoids the
connected curve $\gamma$, the two centres on $\gamma$ lie in the same
complementary component and the two indicators agree.

Multiply \eqref{eq:winding-indicator-telescope} by $w(\eta)$ and sum over
$\eta\in\mathscr W_N^0$. Every placed polygon on the right is counted
exactly $N$ times: summing both the anchored shapes and all $N$ horizontal
translations has this multiplicity even when a shape has a horizontal
stabilizer. Dividing by $N$ yields
\[
 P_N\delta_\gamma\le
 \sum_{\eta\in\mathscr W_N}w(\eta)
             \mathbf1_{\{\gamma\not\sim\eta\}}.
\]
Summing over $\gamma$ and using $P_N\asymp N^{-1}$ in
\eqref{eq:winding-pair-mass} proves
\begin{equation}\label{eq:winding-height}
 \frac1{P_N}\sum_{\gamma\in\mathscr W_N^0}
          w(\gamma)\operatorname{ht}(\gamma)
 \le C_{\rm lat}+\frac{Q_N}{P_N^2}\le C N.
\end{equation}

\paragraph{Normalized arc entries.}
For $1\le g<N$, let $\mathscr A_{N,M}(g)$ be the arcs counted by $K_N(g)$
whose occupied rows lie among $-M,\ldots,-1$. For such an arc define the
finite gas with all windings touching it excluded,
\[
 \Xi_{N,M}^{\gamma}(\nu)=
 \sum_{\substack{\mathcal B\subset\mathscr W_{N,M}\text{ pairwise compatible}\\
                  \eta\cap\gamma=\varnothing\ (\eta\in\mathcal B)}}
       \prod_{\eta\in\mathcal B}\nu w(\eta).
\]
The one-arc entry of $(\mathsf T_\nu(l))^M|0\rangle$, divided by its empty entry,
is exactly
\begin{equation}\label{eq:arc-finite-ratio}
 A_{N,M}(g,\nu)=\sum_{\gamma\in\mathscr A_{N,M}(g)}
       w(\gamma)\frac{\Xi_{N,M}^{\gamma}(\nu)}{\Xi_{N,M}(\nu)}.
\end{equation}
Both numerator and denominator count the same finite winding gas; the
numerator retains precisely the gas configurations compatible with the
specified arc. Put
\[
 W_{N,M}(\gamma)=\sum_{\eta\in\mathscr W_{N,M}}
       w(\eta)\mathbf1_{\{\eta\cap\gamma\ne\varnothing\}}.
\]
This is again bare winding mass.

We recall why all order-$k$ coefficients of the ratios in
\eqref{eq:arc-finite-ratio} have the same sign. For a graph $I$ on $r$
labelled vertices, put
\[
 \phi(I)=\sum_{\substack{G\subset I\\G\text{ connected and spanning}}}
             (-1)^{|E(G)|}.
\]
The connected-graph expansion for a finite hard-core gas gives
\[
 [\nu^r]\log\Xi_{N,M}
 =\frac1{r!}\sum_{\eta_1,\ldots,\eta_r\in\mathscr W_{N,M}}
      \prod_i w(\eta_i)\,\phi(I(\eta_1,\ldots,\eta_r)).
\]
Here $I$ is the incompatibility graph on the labelled indices, so equal
polygons at two indices are joined by an edge. This formula follows by
expanding the pairwise compatibility factors and retaining connected
graphs in the logarithm; see Fern\'andez and Procacci~\cite[Equations~(2.3)--(2.4)]{FernandezProcacci2007}. The value $\phi(I)$ has sign
$(-1)^{r-1}$, or is zero. Indeed the minimum-spanning-tree form of
Penrose's tree-partition identity partitions the connected subgraphs
into intervals with a fixed minimal spanning tree; their alternating
sums are either zero or $(-1)^{r-1}$. See Penrose~\cite{Penrose1967} and
Scott and Sokal~\cite[Lemma~2.2 and Proposition~2.3]{ScottSokal2005}.

Subtracting the logarithm for the unrestricted gas deletes exactly the
tuples with at least one polygon touching $\gamma$. Thus, as formal power
series,
\[
 \begin{gathered}
 \log\frac{\Xi_{N,M}^{\gamma}(\nu)}{\Xi_{N,M}(\nu)}
       =\sum_{r\ge1}(-1)^r b_{r,N,M}(\gamma)\nu^r,\\
 b_{r,N,M}(\gamma)\ge0,\qquad
 b_{1,N,M}(\gamma)=W_{N,M}(\gamma).
 \end{gathered}
\]
After substituting $\nu=-z$, exponentiation has nonnegative
coefficients. Its order-$k$ coefficient contains the term
$W_{N,M}(\gamma)^k/k!$. Therefore
\begin{equation}\label{eq:arc-finite-moment-lower}
 (-1)^k[\nu^k]A_{N,M}(g,\nu)
 \ge\frac1{k!}\sum_{\gamma\in\mathscr A_{N,M}(g)}
                  w(\gamma)W_{N,M}(\gamma)^k\ge0.
\end{equation}

For fixed $N$, the analytic gap makes $A_{N,M}(g,\nu)$ converge
coefficientwise to the normalized eigenvector entry $A_N(g,\nu)$ for
one arc of span $g$. At $\nu=0$ this entry is $K_N(g)$. On the right of
\eqref{eq:arc-finite-moment-lower}, both the set of arcs and the hitting
mass increase with $M$. Write $\mathscr A_N(g)=\bigcup_M\mathscr A_{N,M}(g)$
and $W_N(\gamma)=\lim_M W_{N,M}(\gamma)$. Taking this fixed-$N$ limit
first gives
\begin{equation}\label{eq:arc-infinite-moment-lower}
 (-1)^k[\nu^k]A_N(g,\nu)
 \ge\frac1{k!}\sum_{\gamma\in\mathscr A_N(g)}
                   w(\gamma)W_N(\gamma)^k\ge0.
\end{equation}

An arc from the upper cut to its deepest point must meet every winding
polygon placed strictly between that cut and that point. For a shape
$\eta\in\mathscr W_N^0$, there are at least
$(\operatorname{dep}(\gamma)-\operatorname{ht}(\eta)-C_{\rm lat})_+$
such row placements. Jensen's inequality for the convex positive-part
function and \eqref{eq:winding-height} imply
\begin{equation}\label{eq:arc-hit-depth}
 \begin{split}
 W_N(\gamma)
 &\ge\sum_{\eta\in\mathscr W_N^0}w(\eta)
       (\operatorname{dep}(\gamma)-\operatorname{ht}(\eta)-C_{\rm lat})_+\\
 &\ge P_N(\operatorname{dep}(\gamma)-C_*N)_+
  \ge c_*(\operatorname{dep}(\gamma)/N-C_*)_+,
 \end{split}
\end{equation}
for absolute positive $C_*,c_*$, since $NP_N$ is bounded away from zero.

It remains to connect these coefficients to the derivatives already
bounded in Lemma~\ref{lem:v-higher-twist-jets}. In the empty-row formula
there, the pure $h^N$ term contributes only at order $\nu$. Among the
interval arrangements, only those with one interval survive
differentiation at $h=1$.
Each gap span has $N$ starting positions and interval factor
$C_0(1-h)h^s$, with $C_0=2\cos(3l)$. Hence the exact coefficient identity
for $k\ge1$ is
\[
 f'_{N,k+1}(1)=-C_0N\sum_{g=1}^{N-1}[\nu^k]A_N(g,\nu).
\]
All summands have sign $(-1)^k$ by
\eqref{eq:arc-infinite-moment-lower}. The derivative bound therefore gives
\[
 \sum_{g=1}^{N-1}(-1)^k[\nu^k]A_N(g,\nu)
    =\frac{|f'_{N,k+1}(1)|}{C_0N}\le C_kN^{-1/4}.
\]
Combining this with \eqref{eq:arc-infinite-moment-lower} and
\eqref{eq:arc-hit-depth} proves the required fixed depth moments:
\begin{equation}\label{eq:arc-depth-moments}
 \sum_{g=1}^{N-1}\sum_{\gamma\in\mathscr A_N(g)}w(\gamma)
       \bigl((\operatorname{dep}(\gamma)/N-C_*)_+\bigr)^k
       \le C_kN^{-1/4},\qquad k\ge1.
\end{equation}
The constants are independent of $N$ for each fixed $k$. We will also use
the weaker full moment bound
\begin{equation}\label{eq:arc-full-moments}
 \sum_{g=1}^{N-1}\sum_{\gamma\in\mathscr A_N(g)}w(\gamma)
              (1+\operatorname{dep}(\gamma)/N)^k\le C_k,
 \qquad k\ge0.
\end{equation}
Indeed the unweighted sum is at most $1/C_0$ by
Proposition~\ref{prop:cylinder}, and the part above $C_*N$ is bounded by
\eqref{eq:arc-depth-moments}.

\paragraph{Contractible polygons and the outer arc.}
Let $\mathscr C_N^0$ be the simple contractible cylinder polygons with
top row $0$, each counted once without orientation. Lift such a polygon
to the plane and remove its highest row. That row contains disjoint
intervals pairing consecutive endpoints, say
$a_1<b_1<a_2<b_2<\cdots<a_r<b_r$. Below it, the remaining arcs form a
noncrossing matching, and their union with the intervals is one cycle.
Contract the intervals to ordered boundary points. A noncrossing cycle
through these points must visit them in cyclic order. Thus the lower
matching consists of the consecutive inner arcs $b_i a_{i+1}$ and one
outer arc $a_1b_r$. The outer arc encloses all the inner arcs. Its span
$\ell$ is less than $N$, since otherwise it would meet a horizontal
translate by interlacing of the endpoints. Its depth bounds the depths
of the inner arcs, and
\[
 \operatorname{ht}(P)\le\operatorname{dep}(\text{outer arc})+C_{\rm lat}.
\]

Here is the bound on all possible inner arcs at a prescribed outer
span. At the physical row $h=x^2$, an interval of positive length $s$
has weight $C_0(1-h)h^s$. Put
\[
 a_s=(1-h)h^s\quad(s\ge1),\qquad
 b_g=C_0K_N(g)\quad(1\le g<N),\qquad b_g=0\quad(g\ge N).
\]
Then $\sum_s a_s=h\le1$ and $\sum_g b_g=1-D_N\le1$. After assigning one
factor $C_0$ to each inner gap, the remaining sum for span $\ell$ is
\[
 \mathcal R_N(\ell)=
 \sum_{r\ge1}\ \sum_{\substack{s_i,g_i\ge1\\
                     \sum_i s_i+\sum_i g_i=\ell}}
       \prod_{i=1}^r a_{s_i}\prod_{i=1}^{r-1}b_{g_i}\le1.
\]
To see the inequality, regard $a$ and $b$ as subprobability laws for
alternating positive increments, with any missing mass sent to a
cemetery state. The events that the running sum first equals the fixed
$\ell$ just after the $r$th $a$-increment are disjoint as $r$ varies,
because every increment is positive. Their probabilities are the
displayed summands.

The starting port of the outer arc has $N$ choices. Ignoring mutual
avoidance of the inner arcs only increases their sum; their product
weights are bounded by the renewal sum above. One unassigned interval
factor $C_0$ remains. Therefore
\[
 \begin{split}
 U_{N,k}(\ell)&:=
       \sum_{\alpha\in\mathscr A_N(\ell)}w(\alpha)
          \left(1+\frac{\operatorname{dep}(\alpha)+C_{\rm lat}}N\right)^k,\\
 \sum_{\ell=1}^{N-1}U_{N,k}(\ell)&\le C_k
 \end{split}
\]
by \eqref{eq:arc-full-moments}, and the contractible sum satisfies
\begin{equation}\label{eq:contractible-height-moments}
 \begin{split}
 \sum_{P\in\mathscr C_N^0}w(P)(1+\operatorname{ht}(P)/N)^k
 &\le C_0N\sum_{\ell=1}^{N-1}\mathcal R_N(\ell)U_{N,k}(\ell)\\
 &\le C_kN,\qquad k\ge0.
 \end{split}
\end{equation}
This is the height estimate needed for the planar projection.

\paragraph{Tip tails and a rotated period.}
Fix a seam tip $v$. For a shape in $\mathscr C_N^0$, at most
$\operatorname{ht}(P)+C_{\rm lat}$ row translations can have its disk
contain $v$. The moment of order $k+1$ in
\eqref{eq:contractible-height-moments} consequently gives, for every
integer $k\ge1$ and $L\ge1$,
\begin{equation}\label{eq:tip-height-tail}
 \begin{aligned}
 &\sum_{\substack{P\text{ contractible, enclosing }v\\
                  \operatorname{ht}(P)>L}}w(P)\\
 &\quad\le \sum_{P\in\mathscr C_N^0}w(P)
      (\operatorname{ht}(P)+C_{\rm lat})
         \mathbf1_{\{\operatorname{ht}(P)>L\}}\\
 &\quad\le L^{-k}\sum_{P\in\mathscr C_N^0}w(P)
              (\operatorname{ht}(P)+C_{\rm lat})^{k+1}
 \le C_kN^{k+2}L^{-k}.
 \end{aligned}
\end{equation}
In particular the bare mass with height greater than $N^{1+\epsilon}$
is at most $C_kN^{2-k\epsilon}$.

For clarity, a contractible cylinder polygon enclosing $v$ has a unique
planar lift whose disk contains a specified lift of $v$. Indeed the
boundaries of two translates of a lifted disk are disjoint. Their disks
must then be disjoint or nested; strict nesting is impossible for two
translates of the same bounded disk, since translation preserves area.
Thus their interiors are disjoint. We call the lift containing the
specified tip the centred lift and denote it by $\widetilde P$.

Fix $\epsilon>0$ and put $M=\lfloor N^{1+2\epsilon}\rfloor$. Consider
centred lifts with old height at most $N^{1+\epsilon}$ and diameter
greater than $N^{1+3\epsilon}$. The translation $M e^{i\pi/3}$ changes
the old row coordinate by $M$. For large $N$ this exceeds the old
height, so the lift and every nonzero translate by this vector have
disjoint old row ranges. It therefore projects to a simple contractible
polygon on the cylinder of period $M$ in direction $\pi/3$.

Write $\operatorname{ht}_0$ and $\operatorname{ht}_{\pi/3}$ for occupied-row
heights in the original and rotated orientations. The two row
coordinates are independent linear coordinates in the plane. There is
an absolute $C$ such that
\[
 \diam(\widetilde P)\le C\bigl(\operatorname{ht}_0(\widetilde P)
                         +\operatorname{ht}_{\pi/3}(\widetilde P)\bigr).
\]
The assumed diameter and old height thus force rotated height at least
$c_1N^{1+3\epsilon}$ for large $N$. This projection is injective on the
centred lifts: if two gave the same rotated-cylinder polygon, they would
differ by a translate of its disk, while both disks contain the same
specified tip lift. The same disjoint-disk argument on the rotated
cylinder forces that translate to be zero. Applying
\eqref{eq:tip-height-tail} after the $\pi/3$ lattice
rotation gives
\[
 \begin{aligned}
 &\sum_{\substack{P\text{ enclosing }v,\ \operatorname{ht}_0(\widetilde P)\le N^{1+\epsilon}\\
                  \diam(\widetilde P)>N^{1+3\epsilon}}}w(P)\\
 &\quad\le C_k M^{k+2}(c_1N^{1+3\epsilon})^{-k}
 \le C_kN^{2+4\epsilon-k\epsilon}.
 \end{aligned}
\]
Together with the old-height tail, this proves
\begin{equation}\label{eq:tip-diameter-tail}
 \begin{split}
 q_{N,\epsilon}&:=
 \sup_v\sum_{\substack{P\text{ contractible, enclosing }v\\
                        \diam(\widetilde P)>N^{1+3\epsilon}}}w(P)\\
 &\le C_k\bigl(N^{2-k\epsilon}+N^{2+4\epsilon-k\epsilon}\bigr)=o(1),
 \end{split}
\end{equation}
where one fixes $k>(2+4\epsilon)/\epsilon$ and then lets $N\to\infty$.
This is a bound on total bare activity, not merely on the probability
of a large polygon in a gas.

\subsection{The two seam tips and the planar comparison}

The height estimates permit the projection to the plane. We first
identify the cap matrix element as a positive two-tip partition and
control the cost of horizontal slab walls.

At homogeneous sites let $G_H=\Lambda(T^1(l))^H\Psi$, where the
superscript $1$ denotes trivial spin twist. The cap regularity argument
in Section~\ref{sec:physical-amplitude} represents the lower and upper
caps as limits of empty boundaries extended by physical rows with twists
$i$ and $-i$, respectively; the upper construction uses transpose,
reversal, and spin negation. In a finite such approximation the seam
twist is $i$ below the tip at cut $0$, $1$ between cuts $0$ and $H$, and
$-i$ above the tip at $H$, as indicated in Figure~\ref{fig:seam}.

For an oriented loop let $c_t$ be its signed crossing number of the
upward seam ray from tip $t$, with left-going auxiliary flow positive.
For a counterclockwise contractible loop, $c_t$ is one if its disk
contains $t$ and zero otherwise. Its turn phase is $i$, and its total
full-seam crossing number is zero. For a winding loop oriented to the
left, $c_t$ is one if $t$ lies in its lower complementary component and
zero otherwise; its turn phase is one and its full-seam phase for the
constant twist $i$ is $i$. In either case the two successive twist
ratios are $1/i=(-i)/1=-i$. The sum of the two orientations is therefore
\[
 2\Re\bigl(i(-i)^{c_0+c_H}\bigr)
 =\begin{cases}
   2,&c_0+c_H=1,\\
   0,&c_0+c_H=0\text{ or }2.
  \end{cases}
\]
Thus precisely the contractible loops enclosing one tip and the
windings separating the tips survive, each with activity $2x^{|P|}$.
Every finite cap approximation is the corresponding positive gas in
its finite cylinder. Exhaustion proves $G_H=\mathcal G_N(H)$.

Let $\mathfrak b_N$ be the total bare weight of windings meeting a fixed
horizontal cut. Counting their possible row placements and using
\eqref{eq:winding-height} gives the uniform wall bound
\begin{equation}\label{eq:winding-wall-mass}
 \mathfrak b_N\le\sum_{\eta\in\mathscr W_N^0}w(\eta)
              (\operatorname{ht}(\eta)+C_{\rm lat})
       \le CNP_N\le C.
\end{equation}
The winding-only slab partition in the introduction is
$L_N(h)=\Xi_{N,h}(2)$, after translating the slab. Stacking slabs and
retaining only polygons internal to each gives the lower inequality
below. For the upper inequality, keep the exclusions within the two
slabs and discard all exclusions involving polygons that meet their
common wall. Those polygons have total activity at most
$2\mathfrak b_N$, so their unrestricted gas is at most
$\prod_\eta(1+2w(\eta))\le e^{2\mathfrak b_N}$. Hence
\begin{equation}\label{eq:slab-wall-inequalities}
 L_N(h)L_N(j)\le L_N(h+j)
       \le e^{2\mathfrak b_N}L_N(h)L_N(j).
\end{equation}
The positive annular transfer at fugacity $2$ is primitive by
Section~\ref{sec:physical-amplitude}. Its dominant eigenvalue
$\lambda_N$ is therefore the exponential growth rate of its
empty-to-empty entry $L_N(h)$. Iterate
\eqref{eq:slab-wall-inequalities} on multiples of $h$ and take the
growth rate to obtain
\begin{equation}\label{eq:slab-growth-bounds}
 e^{-2\mathfrak b_N}\lambda_N^h\le L_N(h)\le\lambda_N^h,
 \qquad 0\le\log\lambda_N\le 2P_N.
\end{equation}
For the last inequality, discarding all exclusions gives
$\log L_N(h)\le2\sum_{\eta\in\mathscr W_{N,h}}w(\eta)\le2hP_N$.
In particular, for $H>2N$,
\begin{equation}\label{eq:inward-wall-cost}
 \frac{L_N(H-2N)}{\lambda_N^H}
 \ge e^{-2\mathfrak b_N}\lambda_N^{-2N}
 \ge e^{-2\mathfrak b_N-4NP_N}\ge e^{-C}.
\end{equation}

Write $Z(r)=\mathcal Z_v(r)$ for the planar nesting partition at either
tip. Choose a small absolute $a_0>0$ so that a disk of radius $a_0N$
projects injectively to the period-$N$ cylinder and, when centred at
the lower tip, lies below the cut at $N$; the analogous disk at the
upper tip lies above the cut at $H-N$. For large $H$ the two disks are
disjoint. Their nests coexist with every winding configuration internal
to the reduced slab $[N,H-N]$, and all those windings separate the tips.
Thus
\[
 \mathcal G_N(H)\ge Z(a_0N)^2L_N(H-2N).
\]

For the upper bound use the original slab $[0,H]$. A separating winding
not internal to this slab must meet its lower or upper wall, so such
windings have total activity at most $4\mathfrak b_N$. A contractible
polygon with centred lift of diameter at most
$D=N^{1+3\epsilon}$ lies within distance $D$ of its enclosed tip: the
tip lies in the convex hull of the polygon. The centred lifts at each
tip form a planar nest, whose partition is bounded by $Z(D)$. Finally,
\eqref{eq:tip-diameter-tail} bounds the total activity of the remaining
contractible polygons at both tips by $4q_{N,\epsilon}$. Retain the
exclusions within each of the two small nests and within the
slab-internal winding gas, and discard all other exclusions. In finite
exhaustions this gives, and hence in the limit gives,
\[
 \mathcal G_N(H)\le
 L_N(H)Z(N^{1+3\epsilon})^2
           \exp(4\mathfrak b_N+4q_{N,\epsilon}).
\]
These bounds in particular make the positive full-cylinder sum finite.
Equations \eqref{eq:slab-growth-bounds}--\eqref{eq:inward-wall-cost} now
give the comparison
\begin{equation}\label{eq:planar-amplitude-squeeze}
 e^{-C}Z(a_0N)^2\le
 \lim_{H\to\infty}\lambda_N^{-H}\mathcal G_N(H)
 \le e^{C_\epsilon}Z(N^{1+3\epsilon})^2.
\end{equation}
The limit exists by the physical cap expansion and equals
$c_{\rm phys}(0)=N^{1/6+o(1)}$. The left inequality in
\eqref{eq:planar-amplitude-squeeze}, with monotonicity between successive
$a_0N$, gives
$\limsup_{r\to\infty}\log Z(r)/\log r\le1/12$. The right inequality,
with $N=\lfloor r^{1/(1+3\epsilon)}\rfloor$, gives
$\liminf_{r\to\infty}\log Z(r)/\log r\ge1/(12(1+3\epsilon))$.
Letting $\epsilon\downarrow0$ proves
$Z(r)=r^{1/12+o(1)}$ and Theorem~\ref{thm:amplitude}.
